%% ma: L12-B19-0029
%% lop: 12
%% chuong: 6
%% bai: 19
%% dang: 12.19.5
%% loai: TN4
%% muc: VD
%% dap_an: "A"
%% nguon: "(NBV)-10 ĐỀ KIỂM TRA HK2-TOÁN 12-DUNG CHUNG 3 SÁCH.pdf (H02, Đề số 2), Phần 1 Câu 12"
%% kien_thuc: "Bài 19 (công thức xác suất toàn phần và công thức Bayes)"
%% trang_thai: chinh-thuc
%% ghi_chu: "Bỏ ảnh minh họa hai người đeo khẩu trang (ảnh trang trí)."
\begin{ex}
Ở một địa phương $X$, xác suất để một người lớn trên $40$ tuổi mắc bệnh ung thư là $0{,}05$. Xác suất bác sĩ chẩn đoán đúng một người mắc bệnh ung thư là $0{,}78$ và chẩn đoán sai (không bị ung thư nhưng lại chẩn đoán mắc bệnh) là $0{,}06$. Xác suất để một người trên $40$ tuổi thực sự mắc bệnh ung thư khi nhận được kết quả chẩn đoán bị ung thư bằng
\choice{\True $0{,}40625$}{$0{,}49615$}{$0{,}53225$}{$0{,}62634$}
\loigiai{Gọi $A$: ``người đó thực sự bị ung thư'', $B$: ``người đó được chẩn đoán bị ung thư''. Ta có $P(A)=0{,}05$, $P(\overline A)=0{,}95$, $P(B\mid A)=0{,}78$, $P(B\mid\overline A)=0{,}06$.

Xác suất toàn phần: $P(B)=0{,}05\cdot0{,}78+0{,}95\cdot0{,}06=0{,}096$.

Công thức Bayes: $P(A\mid B)=\dfrac{0{,}05\cdot0{,}78}{0{,}096}=0{,}40625$. Chọn A.}
\end{ex}
